\documentclass[11pt]{article}
\usepackage[utf8]{inputenc}
\usepackage[T1]{fontenc}
\usepackage{lmodern}
\usepackage[margin=1in]{geometry}
\usepackage{csquotes}
\usepackage[style=authoryear,backend=biber,sorting=nyt,maxbibnames=99,giveninits=true,uniquename=init]{biblatex}
\usepackage[hidelinks]{hyperref}
\addbibresource{references.bib}
\begin{document}
\section*{Tile Drainage Research Gaps and Modeling Needs}

\subsection{Overview}

Despite decades of field and modelling work, the literature synthesised here converges on a shared diagnosis: our capacity to predict how subsurface (tile) drainage alters water, sediment, and nutrient export remains limited by sparse field verification, incomplete knowledge of where and how drains are installed, and process representations that are too coarse for the scales at which decisions are made. The recurring research frontier therefore spans four interlocking gaps. First, watershed-scale hydrologic effects of tiling must be assessed with both models and field observations, across events of all sizes and multiple landscapes, rather than inferred from single sites \parencite{Schilling2008,Sloan2017,Sloan2016}. Second, the spatial location, extent, and specifications of drainage networks are largely unknown, motivating calls for mapping, databases, and validated prediction of tile lines \parencite{Naz2009,Zhou2025,Valayamkunnath2020}. Third, existing agricultural models (APEX, DRAINMOD, SWAT, MIKE~SHE, and others) require improved tile, water-table, macropore, and phosphorus routines, together with careful treatment of parameter uncertainty when drainage specifications are unknown at calibration \parencite{Williams2016,Bauwe2016,Guo2018b,Nelson2018}. Fourth, key transport processes---preferential and macropore flow, phosphorus source partitioning, sediment storage and exhaustion, and cold-season dynamics---remain underquantified and must be constrained with longer-term, higher-resolution monitoring \parencite{King2015,Grangeon2021,Trentman2020}. The subsections below organise the identified gaps, hypotheses, and methodological needs around these themes.

\subsection{Watershed-Scale Hydrology and Cross-Scale Integration}

A dominant line of argument is that the hydrologic effects of tile drainage remain inadequately characterised at the watershed scale and must be pursued through combined modelling and field verification. \textcite{Schilling2008} argue that more in-depth research is needed to assess tile drainage effects on watershed hydrology using both models and field data, with specific attention to streamflow recessions and how those recessions differ seasonally between tiled and non-tiled watersheds. Extending this, \textcite{Schilling2012} call for work that quantifies the specific hydrologic effects beneath patterned tiling systems at the watershed scale. \textcite{Sloan2017} frame the outstanding questions in terms of scope, recommending that the influence of tile drains be explored for events of all sizes and at all scales, alongside more comprehensive analyses across different landscapes.

Integrating field-scale responses into basin behaviour is itself an open problem. \textcite{Sloan2016} conclude that hydrologic models must include routing between spatially variable depressional stores, and that a detailed modelling study is needed to understand flow accumulation and attenuation when responses from multiple tiled fields are integrated across a river network and their combined role in large-scale watershed flooding. Related scaling questions concern how the collective drainage network shapes flood behaviour: \textcite{Fonley2021} recommend a comprehensive study of the statistical improvement in flood prediction offered by their model, further study of how the flow duration curve evolves as the tile population evolves, and note that such work could ultimately inform recommendations for where artificial drainage is applied on the landscape. At the regional scale, \textcite{BolandBrien2014} identify a need to decouple the effects of row-crop cover from artificial drainage and to separate the contributions of surface versus subsurface drainage in altering hydrologic response; they hypothesise that bypass-flow networks contribute to hydrologic homogenisation and argue that simple models for tiled landscapes should be developed, in part because long-term data are sparse.

Several studies emphasise interactions between drains and the receiving channel network as a frontier. \textcite{Grangeon2021} suggest future investigation of the interactions between rivers/ditches and tile-drain flows, while \textcite{Clement2017} call for research on the fate and transport of tile-drain effluent after it reaches the surface drainage system. \textcite{Zhou2025} argue for field-scale monitoring of the direct impact of tile flow on downstream water discharge, stream power, and channel geomorphology across multiple locations within watersheds to provide higher-resolution ground-truth measurements. Efforts to extend field-scale drainage models to catchment and watershed scales are ongoing, as with the application of MIKE~SHE to Midwestern tile-drained lands \parencite{Zhou2013}.

\subsection{Locating, Mapping, and Cataloguing Drainage Networks}

A cross-cutting obstacle is that the physical drainage network is often unknown, which propagates uncertainty into every downstream model. At the field scale, \textcite{Amado2018} note that having actual tile locations, and measured tile-flow discharge at the site, would improve model accuracy and calibration, and that additional field data and delineation of the subsurface drainage network are needed to refine the model (alongside a grid-sensitivity study). Predictive mapping of tile lines from imagery is promising but immature: \textcite{Naz2009} call for further validation of predicted tile maps against observed locations from GPS survey, and for methods that minimise the misidentification of image noise as tile lines and that connect broken segments erroneously split during detection. \textcite{Schilling2019} likewise identify improving the ability to locate field tile systems in agricultural watersheds as an ongoing research focus.

At larger scales, \textcite{Valayamkunnath2020} anticipate updating their tile-drainage map repository as new data become available. \textcite{Zhou2025} set out the most explicit data agenda: future studies should report information on subsurface tiles beyond areal extent---particularly tile diameter, depth, and spacing, as well as the mains that collect water from pattern tiles---and a tile-drainage database should be constructed through joint efforts of government agencies, researchers, private landowners, and drainage contractors, with pipe networks mapped using remote sensing and ground-based surveys where possible.

\subsection{Model Structure, Routines, and Parameterisation}

\subsubsection{Improved and New Tile, Water-Table, and Crop Routines}

Many authors call for better process routines within existing agricultural models. \textcite{Williams2016} argue that incorporating identified flow pathways and processes into models such as APEX, DRAINMOD, and SWAT is critical for improving estimates of phosphorus transport in tile-drained landscapes. Within SWAT and its variants, \textcite{Schilling2019} recommend testing new tile-drainage algorithms that give a more physically based representation of tile function and can directly account for pattern-tile configurations; \textcite{Guo2018b} see an opportunity to improve representation of tile systems---especially individual tiles---and functionality at watershed scales, including improving water-table-depth calculation during long dry periods and testing improved routines and curve-number methods on additional individual tiles and watersheds. \textcite{Moriasi2012} call for follow-up studies comparing dynamic and constant surface-depression storage ($S_d$) approaches for tile-flow and streamflow simulation, and for additional testing of the Modified SWAT's tile-flow and water-quality simulation against measured tile flow and tile-flow water-quality data. In APEX, \textcite{Nelson2018} note that improvements in tile-drainage, water-table, and crop-growth routines may be needed, and that new routines for water-table fluctuation and subsurface tile drainage were in development but not yet implemented. For rangeland modelling, \textcite{Chen2016} conclude that a tile-drainage module is needed in the GPFARM-Range model to predict drainage from subsurface-drained pastureland on shorter time scales and in a timely manner. \textcite{Boles2015} report that their new routine's parameter values provide a basis for fine-tuning tile parameterisation, suggest that modellers vary LATKSATF on a hydrologic-response-unit rather than watershed basis, and identify additional studies to build understanding of the routine's effects.

\subsubsection{Parameter Uncertainty When Drainage Specifications Are Unknown}

A distinct methodological caution concerns calibration under unknown drainage specifications. \textcite{Bauwe2016} stress that much care should be taken when calibrating the depth-to-impervious-layer parameter (DEP\_IMP) because it is critical in dividing subsurface flow into tile flow and groundwater flow, that the choice of infiltration model matters when flow components are relevant, and that this is especially important when drainage specifications are unknown and tile-drain parameters are used for calibration. Consistent with this, \textcite{Guo2018b} propose separating DEP\_IMP into two parameters---depth to impervious layer and a soil-permeability factor---to improve representation of the tile system.

\subsubsection{Process Formulations and Higher-Dimensional Models}

Some gaps concern the underlying physics of flow and solute movement. \textcite{Qi2018} suggest that the method of simulating subsurface water flow in the soil matrix should be revisited or replaced by a soil-matrix-potential-based Richards equation, and that the algorithm for particulate-phosphorus loss to tile drainage should be revisited. \textcite{Mohanty1998} argue that tile flow and conservative solute transport would perhaps be better interpreted and modelled with a three-dimensional model considering regional water and chemical management, and that more detailed consideration of regional hydrology and adjacent-bench irrigation/fertiliser scheduling is needed to improve prediction. \textcite{DeSchepper2015} report developing a larger-scale model to simulate seasonal variation in tile-drainage outflow with heterogeneities of clayey tills represented in greater detail, and argue that more realistic simulated water-flow patterns in heterogeneous drained fields are a cornerstone of future nutrient-transport modelling; they anticipate future work focusing on nutrient-transport modelling. \textcite{Garrison1999} propose that integrating spatial images of surface wetness (before emergence) and within-season vegetative index, coupled with final yield, will refine tile-model parameters, and that their tile-water and nitrogen-flow methods may found improved spatial representation of soil-water and nitrogen dynamics.

\subsubsection{Field Verification of Model Predictions}

Calls for direct measurements to validate models recur throughout. \textcite{Muma2017} recommend further validating the CATHY model by installing an automatic tile-drain system for flow measurement, using field measurements of evapotranspiration to validate simulated actual evapotranspiration, and installing a piezometer at the centre of the tile-drain plot to evaluate model performance against water-table elevation. \textcite{Nazari2021} suggest their approach can help inform numerical-model evaluations and assess environmental gradients across sites, and that accounting for redox and other conditions affecting sorption--desorption dynamics between events could reduce uncertainty and improve model fit (Nash--Sutcliffe efficiency).

\subsection{Preferential and Macropore Flow Processes}

Preferential and macropore flow are repeatedly named as critical, underquantified controls. \textcite{Boles2015} identify macropore-related processes as a critical area of future research. \textcite{Williams2016} call for investigation of macropore-flow contributions to tile drains during different periods of the year to clarify their role in annual water budgets, and for work on the interaction between macropores and the surrounding matrix. \textcite{King2015} frame the central questions as quantifying the development, extent, and behaviour of preferential-flow paths under different soils, management practices, and antecedent conditions, and determining what fraction of phosphorus is delivered to tile drains via preferential flow versus matrix flow. \textcite{Nazari2021} recommend that their methodology be evaluated at sites where matrix--macropore exchange forms a larger share of the tile water budget, and that it be considered in water-quality research on practices---such as drainage water management---that directly affect water-table dynamics. Methodologically, \textcite{Nielsen2015} note their staining technique should be most useful in studies of preferential flow and transport in drained soils (for example, leaching of strongly sorbing contaminants), and that longer infiltration times could increase the number of stained pathways reaching the tile lines. \textcite{Ruark2009} argue that further explanation of preferential-flow interpretation requires detailed analyses of hydrograph and chemograph behaviour, and \textcite{BACKLUND1995} call for research on contamination transported by old and new tile systems through preferential-flow paths and on means of modelling it.

\subsection{Phosphorus and Nutrient Transport}

Improving the representation and understanding of nutrient transport---especially phosphorus---is a persistent priority. \textcite{King2015} call for developing and enhancing phosphorus routines in deterministic models to better represent subsurface drainage and P transport, and for measuring phosphorus losses in surface runoff and subsurface drainage simultaneously to determine partitioning between flow paths and avoid unintended consequences. \textcite{Williams2016} stress the importance of differentiating sources of tile-drained phosphorus (e.g., soil versus fertiliser) to understand how P is transported through macropore paths. \textcite{Qi2018} encourage revisiting the particulate-P loss algorithm and suggest measuring the size distribution of soil particles that act as phosphorus carriers, comparing them with soil-matrix pore sizes to enhance understanding of particulate-P transport. \textcite{Geohring2001} note that additional research would help determine which processes most significantly reduce total-phosphorus loadings to tile drains. \textcite{Zimmer2016} recommend more intensive monitoring of discharge and P forms from tile-drained areas, and considering retention of phosphorus at tile-drain outlets to prevent excessive input to water bodies.

Broader nutrient questions extend to nitrogen and receiving-water impacts. \textcite{Clement2017} call for research on P retention and uptake lengths in reaches with and without best management practices, on sampling tile drains from different areas of the watershed, on the potential effects of nitrogen given its prevalence in inorganic and organic fertilisers, and on the role of internal P loading from lake sediments in driving algal blooms. \textcite{King2018} argue that additional work is needed to understand phosphorus losses in subsurface tile-drained landscapes and to mitigate harmful and nuisance algal blooms in receiving waters, examining both individual and stacked nutrient-management practices and the processes controlling P transport across fields with varying characteristics. \textcite{Ruark2009} note that both the nature and drivers of dissolved-organic-carbon dynamics in relation to effluent applications require further investigation, and that the ecological implications of DOC from tile drainage remain largely unknown.

\subsection{Sediment Dynamics and Storage}

Sediment transfer through drains is highlighted as an emerging area requiring longer records and targeted tracers. \textcite{Grangeon2021} recommend monitoring tile-drained catchments over longer periods to gain insight into sediment storage and exhaustion, further quantifying the factors controlling sediment storage from the cultivated topsoil to the drains during ephemeral flow, using short-lived fallout radionuclides or alternative low-cost measurements to quantify sediment storage and transfer, and further quantifying the impacts of these transfers on sediment and contaminant exports from drained catchments. \textcite{Gaillot2023b} argue that the influence of cropping practices on soil-surface characteristics and tile-drain connectivity---and on the temporal dynamics of suspended sediment---should be explored further; they recommend combining high-temporal-resolution acquisition of hydrosedimentary fluxes with the dynamic evolution of soil structure at the surface, in the topsoil, and in the subsoil, and suggest that a suspended-sediment tracer allowing event-scale discrimination could clarify sediment dynamics. At the reach scale, \textcite{Ahiablame2010} call for investigations accounting for manure application, soil types, and nutrient loadings from individual tile outlets to examine the impact of tile inputs on sediments across different reach scales and ditch conditions (vegetated, non-vegetated, dredged, and non-dredged).

\subsection{Management Practices, Trade-offs, and Mitigation}

A further frontier concerns rigorous, holistic evaluation of management and treatment practices. \textcite{King2015} call for holistic evaluations of management and treatment practices---including cumulative effects and trade-offs between surface runoff and subsurface drainage---and for exploring practices that target phosphorus in drainage water, such as tillage directly over tile lines and withholding phosphorus application over tile lines. \textcite{King2016} recommend investigating cover crops, drainage water management, and 4R nutrient management as strategies to reduce nitrogen and phosphorus transport, particularly practices targeting the non-growing season to decrease annual loads. \textcite{Williams2016} call for examining different tillage operations (moldboard plow, chisel plow) and application methods (banding or injection of fertiliser). \textcite{Kokulan2021} recommend assessing tile drainage's role in hydrochemical export together with practices such as fertilisation, tillage, and variable tile configuration to develop a complete regional picture. \textcite{Qi2018} suggest testing further off-season management and modelling practices under manure-phosphorus application and different soil and climate conditions. In the context of agricultural-waste land application, \textcite{BACKLUND1995} call for research on nutrient leaching, interception, denitrification, and fecal-coliform transport; on the interaction between cedar wood chips and iron-ocher buildup (including chemical substitutes or filter-fabric applications); on reuse of drainage water and its impact on crop production and downstream water quality; and, where contamination exceeds acceptable levels, on alternatives such as tile-outflow pumpback systems, reuse, or other outlet options. \textcite{Fortin2002} note that evaluating the effects of different soil treatments on chemical movement to tile drains under natural rainfall probably requires long-term experiments to capture climatic variability.

\subsection{Seasonality, Cold Climates, and Long-Term Monitoring}

Several gaps concern temporal coverage. \textcite{King2015} identify the need to expand research on hydrology and phosphorus transport in cold climates, especially the winter and snowmelt periods that are underrepresented in the literature. \textcite{Trentman2020} recommend prioritising seasonal sampling across a range of antecedent-precipitation conditions in light of a changing climate, note that the effect of snow on tile drainage and nutrient export is a topic of interest, and observe that soluble-reactive-phosphorus reduction estimates could be better constrained with continuous (e.g., daily) tile-drain measurements. \textcite{Millar2025} call for multiyear studies of tile-drain nutrient loads and for coverage of important periods outside the growing season, including winter, postharvest, and early spring. Longer and more intensive observation is also needed to diagnose drain performance itself: \textcite{Rimidis2003} argue that a more elaborate measuring set-up and long-term observations are required to ascertain the cause of drain malfunctioning, recommending additional piezometers close to the drain to measure approach-flow head loss and to distinguish head due to entrance resistance from backpressure within the drain.

\subsection{Data, Tracers, and Methodological Needs}

Finally, the reviewed studies point to methodological and data infrastructure requirements that cut across the themes above. Beyond drainage-network mapping and databases \parencite{Zhou2025,Valayamkunnath2020,Naz2009}, authors call for tracer- and chemistry-based approaches to partition flow components: \textcite{Frederiksen2023} suggest process-based modelling or water-chemistry studies using water isotopes or conservative tracers to estimate fast-flowing components in quick flow and the tile-flow component in baseflow, and recommend adding predictor variables such as landscape characteristics to better explain spatial variation in tile-flow generation and improve modelling accuracy and precision. \textcite{Schilling2019} propose adding less-sensitive parameters such as specific conductance to end-member mixing analysis to better account for seasonal deployment and operational variations. Collectively, these needs---higher-resolution and longer-term monitoring, event-scale tracers, three-dimensional and physically based process models, and open shared datasets on drainage-network specifications---define a research agenda in which improved field verification and improved model structure advance together rather than separately \parencite{Schilling2008,Bauwe2016,Zhou2025}.

\printbibliography
\end{document}
